\section{Limitations}
\label{sec:limits}

The questions are factual and short, and their answers come from a closed set of 62 countries, cities, languages and sports; three of the four kinds are names, and for the one kind that is not (sports) the jump is much smaller and late readability lower, so whether places behave the same way for common nouns in general is open. Places for open-ended answers, multi-step reasoning or long contexts may look different. We studied four instruction-tuned models between 1 and 9 billion parameters, one prompt format each, and greedy answers. Readability uses a whitening of the state; the tightness ratio does not, and it falls fastest at the same depth, but both measure where states sit, not the computation that moves them there. The step fit compares two simple shapes; it shows that a sudden change describes the data much better than a constant contraction, not that the change happens within a single layer. Moves at late layers partly write the answer directly, so the strongest causal evidence is the depth at which moves start to work and the failure of random moves. Wikidata answers are a reference, not ground truth; we removed one unreliable relation--answer pair after a manual check, and a small number of wrong labels surely remain. The rewordings were written by a language model and filtered automatically.
